% Copyright (c) 2026 Benjamin Stanley Frohman. All rights reserved.
% Released under Apache 2.0 license as described in the file LICENSE.
% Authors: Benjamin Stanley Frohman (ORCID: 0009-0006-7068-3718)
\documentclass[11pt]{article}
\usepackage[margin=1in]{geometry}
\usepackage{amsmath,amssymb,amsthm}
\usepackage{booktabs}
\usepackage{hyperref}

\newtheorem{theorem}{Theorem}[section]
\newtheorem{proposition}[theorem]{Proposition}
\newtheorem{definition}[theorem]{Definition}
\newtheorem{remark}[theorem]{Remark}

\title{\textbf{Motion-First Geometry, Lorentz Ether vs.\ Special Relativity, Poincar\'e's Fuchsian/Lorentz Group Theory, and Casimir/Wigner $\mathrm{ISO}(1,3)$ Classification}}
\author{Benjamin Stanley Frohman\\
\small Independent Researcher \quad|\quad ORCID: \href{https://orcid.org/0009-0006-7068-3718}{0009-0006-7068-3718}\\
\small \texttt{frohmanbenjamin@gmail.com}}
\date{Origin Provenance: May 19, 2026 (Google Drive Archive: July 7, 2026, 14:21:41 CDT) \\ \small Companion Physics \& Group Theory Quad-Treatise \quad|\quad Licensed under Apache-2.0}

\begin{document}
\maketitle

\begin{abstract}
This treatise synthesizes four foundational investigations authored in May 2026 during the development of the Frohmanian Symplectic Tether program:
(I) \emph{Motion-First Geometry and Relational Dynamics}, contrasting relational kinematic axioms (Heraclitus, Leibniz, Barbour--Bertotti shape space, and Rovelli relational quantum mechanics) against static substantivalist spacetime blocks;
(II) \emph{Lorentz Ether Theory (LET) vs.\ Einstein Special Relativity (SR)}, analyzing empirical equivalence under the Lorentz transformations alongside ontological divergence between dynamical stress in a preferred frame and Minkowski metric kinematics;
(III) \emph{Poincar\'e's Group-Theoretic Synthesis (1880--1906)}, tracing the mathematical continuity from Fuchsian automorphic groups and the hyperbolic disk model of $\mathrm{PSL}(2,\mathbb{R}) \cong \mathrm{SO}^+(1,2)$ to the 1905/1906 Palermo paper establishing the Lorentz group $\mathrm{SO}^+(1,3)$, its four-vector invariants, and the 10-dimensional inhomogeneous Poincar\'e group $\mathrm{ISO}(1,3) = \mathbb{R}^{1,3} \rtimes \mathrm{SO}^+(1,3)$; and
(IV) \emph{Casimir Invariants and Wigner's 1939 Unitary Irreducible Representation Classification}, deriving how the two polynomial Casimir operators $C_1 = P_\mu P^\mu$ and $C_2 = W_\mu W^\mu$ (with Pauli--Lubanski pseudovector $W^\mu = \tfrac{1}{2}\varepsilon^{\mu\nu\rho\sigma} J_{\nu\rho} P_\sigma$) classify elementary relativistic particles solely by invariant mass $m^2$ and spin/helicity $(s,\lambda)$ across the massive ($m^2>0$, $\mathrm{SU}(2)$ little group), massless ($m^2=0$, $\mathrm{ISO}(2)$ little group), and tachyonic ($m^2<0$, $\mathrm{SO}(1,2)$ little group) orbits.
\end{abstract}

\tableofcontents

\section{Provenance and Archival Record}
\begin{itemize}
  \item \textbf{Author:} Benjamin Stanley Frohman (\texttt{ORCID: 0009-0006-7068-3718}).
  \item \textbf{Origin Provenance Date:} May 19, 2026 (archived in Google Drive as \texttt{Grok Validation-proof.rtf}, File ID \texttt{1P0tCYaDZEiF7riwApIqzfauMzk9vGGqQ}, \texttt{2026-07-07T19:21:41.060Z}, 14:21:41 CDT).
  \item \textbf{Repository Companion Files:} Includes both the verbatim archival RTF/plaintext transcript (\texttt{Grok\_Validation\_Proof\_2026-05-19.rtf}) and this typeset mathematical monograph.
\end{itemize}

\section{Part I: Motion-First Physics and Relational Shape Dynamics}
Standard textbook presentations of mechanics and relativity begin by postulating a fixed background manifold $\mathcal{M}$ equipped with a metric tensor $g_{\mu\nu}$ and subsequently define motion as a parameterized curve $\gamma : \mathbb{R} \to \mathcal{M}$. By contrast, a \emph{motion-first} ontology takes relational flux and state transitions as primitive, recovering spatial geometry and temporal duration as emergent invariants of the relational configuration space.

\subsection{Shape Space and Temporal Emergence}
Let $\mathcal{Q} = (\mathbb{R}^3)^N$ denote the configuration space of $N$ particles, and let $\mathrm{Sim}(3) = \mathbb{R}^3 \rtimes (\mathrm{SO}(3) \times \mathbb{R}_+)$ denote the similarity group of translations, rotations, and global dilatations. In relational mechanics (Leibniz--Mach--Barbour--Bertotti), the physical configuration space is the quotient \emph{shape space}
\[
\mathcal{S} = \mathcal{Q} / \mathrm{Sim}(3).
\]
Rather than evolving with respect to an external Newtonian clock $t$, trajectories correspond to geodesics on $\mathcal{S}$ under a Jacobi-type Maupertuis action
\[
S_{\mathrm{rel}} = \int \sqrt{2(E - V(q))\, T_{\mathrm{best-matched}}(\dot{q},\dot{q})}\,d\lambda,
\]
where the ephemeral time increment $dt_{\mathrm{eph}}$ is defined directly from the relational kinetic change itself:
\[
dt_{\mathrm{eph}} = \sqrt{\frac{T_{\mathrm{best-matched}}(dq,dq)}{2(E - V(q))}}.
\]
Thus time is not an external container in which motion occurs; time is the canonical measure of relational change.

\section{Part II: Lorentz Ether Theory vs.\ Einstein's Special Relativity}
Both Lorentz--Poincar\'e Ether Theory (LET, 1892--1905) and Einstein's Special Relativity (SR, 1905) share the exact same kinematic coordinate transformation law between inertial frames moving at relative velocity $v$ along the $x$-axis:
\[
t' = \gamma\!\left(t - \frac{vx}{c^2}\right), \qquad x' = \gamma(x - vt), \qquad y' = y, \qquad z' = z, \qquad \gamma = \frac{1}{\sqrt{1 - v^2/c^2}}.
\]
Consequently, both frameworks yield identical predictions for all Michelson--Morley, Kennedy--Thorndike, Ives--Stilwell, and synchrotron experiments. Their divergence is structural and geometric:
\begin{center}
\begin{tabular}{p{3.4cm}p{5.5cm}p{5.5cm}}
\toprule
\textbf{Feature} & \textbf{Lorentz Ether Theory (LET)} & \textbf{Einstein Special Relativity (SR)} \\
\midrule
Ontological Primitive & Preferred rest frame $\Sigma_0$ (ether medium) with absolute simultaneity & No preferred inertial frame; Minkowski spacetime $(\mathbb{R}^{1,3},\eta_{\mu\nu})$ \\
Length Contraction & Dynamical electromagnetic compression of intermolecular forces moving through $\Sigma_0$ & Reciprocal geometric projection of worldtubes across non-orthogonal simultaneity hyperplanes \\
Time Dilation & Dynamical slowing of physical oscillators ("local time" $t' = t - vx/c^2$ upgraded to exact $\gamma$-law) & Proper time $\Delta\tau = \int \sqrt{-ds^2}/c$ along timelike curves \\
Group Symmetry & Observational symmetry masking an asymmetric substrate & Fundamental spacetime symmetry group $\mathrm{ISO}(1,3)$ \\
\bottomrule
\end{tabular}
\end{center}

\section{Part III: Poincar\'e's Group-Theoretic Arc (1880--1906)}
Henri Poincar\'e's 1905 Comptes Rendus note and 1906 Rendiconti del Circolo Matematico di Palermo memoir (\emph{Sur la dynamique de l'\'electron}) were the culmination of a 25-year program applying Lie's theory of continuous transformation groups to geometry and mathematical physics:
\begin{enumerate}
  \item \textbf{Fuchsian Groups and Hyperbolic Geometry (1880--1882):} Poincar\'e recognized that discrete Fuchsian groups of fractional linear transformations $z \mapsto \frac{az+b}{cz+d}$ preserving the unit disk $\mathbb{D}$ are isomorphic to isometries of the hyperbolic plane, with $\mathrm{PSL}(2,\mathbb{R}) \cong \mathrm{SO}^+(1,2)$ preserving the ternary quadratic form $x_0^2 - x_1^2 - x_2^2$.
  \item \textbf{Foundations of Geometry via Lie Groups (1887--1898):} In \emph{On the Fundamental Hypotheses of Geometry} (1887) and \emph{Science and Hypothesis} (1902), Poincar\'e demonstrated that homogeneous two-dimensional Riemannian geometries of constant curvature correspond bijectively to three-parameter Lie groups ($\mathrm{SO}(3)$, $\mathrm{ISO}(2)$, $\mathrm{SO}^+(1,2)$), formulating the thesis that geometry is the invariant theory of a transformation group arising from physical displacements.
  \item \textbf{The Lorentz Group and Four-Vector Formalism (1905--1906):} Poincar\'e proved that the transformations preserving Maxwell's equations form a Lie group (which he named the \emph{Lorentz group} $\mathrm{O}(1,3)$), introduced the imaginary time coordinate $x_4 = ict$ so that the interval
  \[
  s^2 = x^2 + y^2 + z^2 - c^2 t^2 = x_1^2 + x_2^2 + x_3^2 + x_4^2
  \]
  is preserved as a four-dimensional orthogonal rotation, constructed four-vectors for position, velocity, current density, and force, and identified the infinitesimal generators of boosts $K_i$ and rotations $J_i$.
\end{enumerate}

\section{Part IV: Casimir Invariants and Wigner's $\mathrm{ISO}(1,3)$ Classification}
The inhomogeneous Lorentz (Poincar\'e) group $\mathcal{P} = \mathrm{ISO}^+(1,3) = \mathbb{R}^{1,3} \rtimes \mathrm{SO}^+(1,3)$ is a 10-dimensional non-compact Lie group generated by four spacetime translations $P^\mu$ and six antisymmetric Lorentz generators $J^{\mu\nu} = -J^{\nu\mu}$ (with $J^i = \tfrac{1}{2}\varepsilon^{ijk}J_{jk}$ and $K^i = J^{0i}$) satisfying the Lie algebra $\mathfrak{iso}(1,3)$:
\begin{align}
[P^\mu, P^\nu] &= 0, \\
[J^{\mu\nu}, P^\rho] &= i\hbar \bigl(\eta^{\nu\rho} P^\mu - \eta^{\mu\rho} P^\nu\bigr), \\
[J^{\mu\nu}, J^{\rho\sigma}] &= i\hbar \bigl(\eta^{\nu\rho} J^{\mu\sigma} - \eta^{\mu\rho} J^{\nu\sigma} - \eta^{\nu\sigma} J^{\mu\rho} + \eta^{\mu\sigma} J^{\nu\rho}\bigr).
\end{align}

\begin{theorem}[Casimir Invariants and Wigner's Little-Group Classification (1939)]
Let $W^\mu := \frac{1}{2}\varepsilon^{\mu\nu\rho\sigma} J_{\nu\rho} P_\sigma$ denote the Pauli--Lubanski pseudovector, which satisfies $P_\mu W^\mu = 0$ and $[P^\mu, W^\nu] = 0$. The center $\mathcal{Z}(\mathcal{U}(\mathfrak{iso}(1,3)))$ of the universal enveloping algebra is generated by the two quadratic and quartic Casimir operators:
\[
C_1 = P_\mu P^\mu, \qquad C_2 = W_\mu W^\mu.
\]
On any unitary irreducible representation $\mathcal{H}$ of the universal cover $\mathbb{R}^{1,3} \rtimes \mathrm{SL}(2,\mathbb{C})$ (by Schur's Lemma):
\begin{enumerate}
  \item \textbf{Massive Class ($C_1 = m^2 c^2 > 0$, $P^0 > 0$):} In the rest frame $p^\mu = (mc, \mathbf{0})$, the stabilizer (little group) is $\mathrm{SU}(2)$ (double cover of $\mathrm{SO}(3)$), $W^\mu = (0, mc\mathbf{S})$, and
  \[
  C_1 = m^2 c^2, \qquad C_2 = -m^2 c^2 \hbar^2 s(s+1), \qquad s \in \bigl\{0, \tfrac{1}{2}, 1, \tfrac{3}{2}, 2, \dots\bigr\},
  \]
  yielding $2s+1$ polarization states $s_z \in \{-s, \dots, +s\}$.
  \item \textbf{Massless Finite-Spin Class ($C_1 = 0$, $C_2 = 0$, $P^\mu \neq 0$):} In the standard lightlike frame $p^\mu = (E/c, 0, 0, E/c)$, the little group is the Euclidean group $\mathrm{ISO}(2) = \mathbb{R}^2 \rtimes \mathrm{SO}(2)$ (double cover). When the non-compact "translational" generators $\Pi_1, \Pi_2$ act trivially, $W^\mu = \lambda P^\mu$ where the helicity $\lambda = \frac{\mathbf{J}\cdot\mathbf{P}}{|\mathbf{P}|} \in \{0, \pm\tfrac{1}{2}, \pm 1, \pm\tfrac{3}{2}, \pm 2, \dots\}$ is Lorentz-invariant.
  \item \textbf{Massless Continuous-Spin Class ($C_1 = 0$, $C_2 = -\rho^2 < 0$):} Non-trivial action of $\mathbb{R}^2 \subset \mathrm{ISO}(2)$ yields infinite-dimensional helicity towers.
  \item \textbf{Tachyonic Spacelike Class ($C_1 = -\mu^2 c^2 < 0$):} Little group $\mathrm{SO}^+(1,2) \cong \mathrm{PSL}(2,\mathbb{R})$, whose non-trivial unitary representations are infinite-dimensional.
\end{enumerate}
\end{theorem}

\section*{License and Copyright}
Copyright \copyright\ 2026 Benjamin Stanley Frohman. All rights reserved. Licensed under the Apache License, Version 2.0 (\url{http://www.apache.org/licenses/LICENSE-2.0}).
\end{document}
